\documentclass{article}

% if you need to pass options to natbib, use, e.g.:
\PassOptionsToPackage{numbers, compress}{natbib}
% before loading neurips_2025

% The authors should use one of these tracks.
% Before accepting by the NeurIPS conference, select one of the options below.
% 0. "default" for submission
 \usepackage[preprint]{neurips_2025}
% the "default" option is equal to the "main" option, which is used for the Main Track with double-blind reviewing.
% 1. "main" option is used for the Main Track
%  \usepackage[main]{neurips_2025}
% 2. "position" option is used for the Position Paper Track
%  \usepackage[position]{neurips_2025}
% 3. "dandb" option is used for the Datasets & Benchmarks Track
 % \usepackage[dandb]{neurips_2025}
% 4. "creativeai" option is used for the Creative AI Track
%  \usepackage[creativeai]{neurips_2025}
% 5. "sglblindworkshop" option is used for the Workshop with single-blind reviewing
 % \usepackage[sglblindworkshop]{neurips_2025}
% 6. "dblblindworkshop" option is used for the Workshop with double-blind reviewing
%  \usepackage[dblblindworkshop]{neurips_2025}

% After being accepted, the authors should add "final" behind the track to compile a camera-ready version.
% 1. Main Track
 % \usepackage[main, final]{neurips_2025}
% 2. Position Paper Track
%  \usepackage[position, final]{neurips_2025}
% 3. Datasets & Benchmarks Track
 % \usepackage[dandb, final]{neurips_2025}
% 4. Creative AI Track
%  \usepackage[creativeai, final]{neurips_2025}
% 5. Workshop with single-blind reviewing
%  \usepackage[sglblindworkshop, final]{neurips_2025}
% 6. Workshop with double-blind reviewing
%  \usepackage[dblblindworkshop, final]{neurips_2025}
% Note. For the workshop paper template, both \title{} and \workshoptitle{} are required, with the former indicating the paper title shown in the title and the latter indicating the workshop title displayed in the footnote.
% For workshops (5., 6.), the authors should add the name of the workshop, "\workshoptitle" command is used to set the workshop title.
% \workshoptitle{WORKSHOP TITLE}

% "preprint" option is used for arXiv or other preprint submissions
 % \usepackage[preprint]{neurips_2025}

% to avoid loading the natbib package, add option nonatbib:
%    \usepackage[nonatbib]{neurips_2025}

\usepackage[utf8]{inputenc} % allow utf-8 input
\usepackage[T1]{fontenc}    % use 8-bit T1 fonts
\usepackage{hyperref}       % hyperlinks
\usepackage{url}            % simple URL typesetting
\usepackage{booktabs}       % professional-quality tables
\usepackage{amsfonts}       % blackboard math symbols
\usepackage{nicefrac}       % compact symbols for 1/2, etc.
\usepackage{microtype}      % microtypography
\usepackage{xcolor}         % colors
\usepackage{amsmath}
\bibliographystyle{unsrtnat}

% Note. For the workshop paper template, both \title{} and \workshoptitle{} are required, with the former indicating the paper title shown in the title and the latter indicating the workshop title displayed in the footnote. 
\title{Few-shot Domain Adaptation for Diffusion Models}


% The \author macro works with any number of authors. There are two commands
% used to separate the names and addresses of multiple authors: \And and \AND.
%
% Using \And between authors leaves it to LaTeX to determine where to break the
% lines. Using \AND forces a line break at that point. So, if LaTeX puts 3 of 4
% authors names on the first line, and the last on the second line, try using
% \AND instead of \And before the third author name.


\author{%
  Yuhang Wu \\
  \texttt{yuw304@ucsd.edu} \\
  \And
  Zhelin Chu \\
  \texttt{z4chu@ucsd.edu} \\
  % examples of more authors
  % \And
  % Coauthor \\
  % Affiliation \\
  % Address \\
  % \texttt{email} \\
  % \AND
  % Coauthor \\
  % Affiliation \\
  % Address \\
  % \texttt{email} \\
  % \And
  % Coauthor \\
  % Affiliation \\
  % Address \\
  % \texttt{email} \\
  % \And
  % Coauthor \\
  % Affiliation \\
  % Address \\
  % \texttt{email} \\
}


\begin{document}


\maketitle


\begin{abstract} Few-shot domain adaptation aims to adapt pre-trained diffusion models to target domains with minimal data. However, standard fine-tuning often leads to overfitting and loss of prior knowledge. This proposal evaluates Low-Rank Adaptation (LoRA) on Stable Diffusion v1.5 for efficient domain shifting. By constraining updates to the UNet's attention layers via low-rank matrices, we aim to achieve high fidelity to target domains while maintaining generative diversity. Our approach is benchmarked against zero-shot and naive fine-tuning baselines to demonstrate its robustness in data-constrained scenarios. \end{abstract}

\section{Introduction}

Few-shot visual domain adaptation addresses the following practical challenge: 
\emph{given a powerful pre-trained text-to-image diffusion model and only 1--5 images from a new visual domain or artistic style, how can we adapt the model to generate images consistent with that domain while still responding to arbitrary text prompts?} 
In this setting, supervision comes solely from a small set of target-domain images (e.g., a specific painting style, medical imaging pattern, proprietary texture, or custom visual aesthetic). The objective is not to learn a new textual concept, but to shift the model's \emph{visual generation distribution} toward the target style while preserving its compositional text-conditioning capability.

This problem is practically important because collecting large-scale domain-specific datasets is often infeasible. Applications such as personalized art generation, niche stylistic rendering, or specialized visual domains require rapid adaptation from only a few examples. A successful few-shot adaptation method must satisfy three criteria simultaneously: 
(i) \textit{fidelity} — generated images should exhibit the target domain's visual characteristics; 
(ii) \textit{diversity} — outputs should not collapse into memorized replicas of the limited training samples; and 
(iii) \textit{prompt controllability} — the model should remain responsive to arbitrary and compositional text prompts.

However, adapting diffusion models under extreme data scarcity is inherently unstable. Naive fine-tuning often overfits rapidly, leading to a collapse where different prompts produce nearly identical outputs resembling the training images. This phenomenon reflects a shift of probability mass toward a few memorized modes and degradation of the original broad visual prior. Prior work on few-shot diffusion adaptation has shown that diffusion models can suffer from diversity collapse and loss of fine-grained texture under limited supervision \cite{zhu2023fewshot}. Meanwhile, full-parameter personalization methods may achieve strong fidelity but require high computational resources and risk degrading general prompt understanding \cite{ruiz2023dreamboothfinetuningtexttoimage}.

In this project, we explore a lightweight and reproducible alternative based on \textbf{Low-Rank Adaptation (LoRA)} \cite{hu2021loralowrankadaptationlarge} applied to \textbf{Stable Diffusion v1.5}. Rather than updating the entire UNet, we insert trainable low-rank matrices into selected attention layers (cross-attention and self-attention) while freezing the remaining parameters. Concretely, for each targeted projection matrix, we learn a low-rank update while keeping the original weight fixed. This constrains adaptation to a low-dimensional subspace of parameter updates, reducing VRAM requirements and limiting drastic deviations from the pre-trained solution.

Conceptually, our hypothesis is that low-rank parameterization acts as an implicit regularizer for visual domain shift. By restricting how the model can change, LoRA may mitigate overfitting while preserving prompt controllability. Unlike objective-level regularization methods that explicitly enforce diversity through additional losses, our approach investigates whether architectural constraints alone are sufficient to stabilize few-shot visual adaptation.

In this project, we make the following contributions. First, we establish a controlled benchmark protocol for extreme few-shot visual domain adaptation (1--5 images), explicitly evaluating fidelity, diversity, and prompt controllability. Second, we systematically analyze how LoRA hyperparameters (rank selection, layer targeting, and training duration) influence overfitting behavior and generative stability. Third, we compare LoRA against zero-shot generation and naive fine-tuning to clarify the trade-off between computational efficiency and adaptation capacity under severe data constraints.

The proposed approach offers several practical advantages: (i) low computational cost and reduced memory usage; (ii) modularity, as LoRA weights can be stored and swapped independently of the base model; and (iii) improved retention of the original generative prior compared to full fine-tuning. Nevertheless, limitations remain. Extremely small and homogeneous datasets may still induce memorization; restricting updates to attention layers may limit expressivity for large structural domain shifts; and evaluation of prompt controllability retains a partially qualitative component. These considerations guide our experimental design and ablation analysis in subsequent sections.


\section{Related Works}
\label{related_works}

Adapting diffusion models to new visual domains has been explored through several approaches, differing in parameter update scale and supervision signals.

\paragraph{Full-parameter fine-tuning.}
Methods such as DreamBooth \cite{ruiz2023dreamboothfinetuningtexttoimage} adapt text-to-image diffusion models by updating large portions of the UNet. While effective in reproducing target appearances, these approaches are computationally expensive and may introduce language drift, reducing generalization to prompts outside the adapted domain.

\paragraph{Embedding-only adaptation.}
Textual Inversion \cite{gal2022imageworthwordpersonalizing} optimizes only a small number of text embeddings while freezing the diffusion backbone. Although computationally efficient, embedding-only methods may lack sufficient representational flexibility to capture complex visual style shifts or texture patterns when supervision is extremely limited.

\paragraph{Parameter-efficient fine-tuning.}
LoRA \cite{hu2021loralowrankadaptationlarge} introduces low-rank weight updates as a parameter-efficient alternative to full fine-tuning. By modifying internal attention projections while keeping the majority of parameters frozen, LoRA balances expressivity and computational efficiency. In visual domain adaptation, this approach implicitly constrains how far the generation distribution can shift.

\paragraph{Few-shot diffusion adaptation via objective regularization.}
Zhu et al.\ \cite{zhu2023fewshot} investigate few-shot adaptation of diffusion models and propose explicit regularization strategies to prevent diversity collapse and high-frequency degradation. Their work emphasizes the importance of preserving relational structure among generated samples. 

In contrast, our approach focuses on \emph{architectural constraints} rather than loss-level regularization. While Zhu et al.\cite{zhu2023fewshot} constrain what the model learns through additional objectives, we constrain how the model can change via low-rank parameterization. This complementary perspective allows us to explore whether structural parameter efficiency alone can stabilize few-shot visual domain adaptation.



% \section{Proposed Methods}
% \label{methods}
% Your proposed solution to solve this problem.

% \textbf{Backbone}: Stable Diffusion v1.5

% \textbf{Baseline}: 
% 1. Zero-shot, 
% 2. Naive few-shot fine-tuning:  The simplest fine-tuning approach involves using the standard diffusion training objective (MSE of noise prediction) on few-shot images.

% \textbf{Main method}: 
% Using LoRA for parameter-efficient fine-tuning:

% LoRA is inserted only into the attention layers (cross-attention/self-attention) of the UNet, while all other parameters are frozen; The training objective still uses the standard diffusion loss, without introducing complex additional networks;
\section{Proposed Methods}

\label{methods}

We formalize few-shot domain adaptation for diffusion models as follows. 
Let $\theta_0$ denote the parameters of a pre-trained Stable Diffusion v1.5 model trained on a large-scale source distribution $\mathcal{D}_s$. 
We are given a small target dataset $\mathcal{D}_t = \{x_i\}_{i=1}^N$ with $N \in [1,5]$, sampled from a target visual domain. 
Our objective is to obtain adapted parameters $\theta$ such that the conditional generation distribution 
$p_\theta(x \mid c)$ shifts toward the target domain while preserving prompt controllability and generative diversity.


\subsection{Datasets}

We evaluate our method on three publicly available datasets spanning natural images, artistic stylization, and medical imagery, enabling controlled few-shot domain adaptation experiments across diverse visual distributions.

\textbf{ImageNet.}
ImageNet is a large-scale natural image dataset containing over 1 million labeled images across 1,000 object categories \cite{deng2009imagenet}. 
Although originally designed for classification, it provides a rich and diverse source distribution. 
We construct few-shot target domains by sampling 1--5 images from selected categories to simulate extreme low-data adaptation scenarios.

\textbf{FS2K (Face Sketch 2000).}
FS2K is a face sketch dataset consisting of approximately 2,000 sketch-style face images \cite{fan2022fs2k}. 
It represents a stylized visual domain distinct from natural photographs, making it suitable for evaluating style-oriented domain shifts (e.g., photo-to-sketch adaptation). 
We sample a small subset of sketch images as the target domain to study few-shot stylistic adaptation.

\textbf{MedMNIST v2.}
MedMNIST v2 is a lightweight collection of biomedical image classification datasets derived from multiple medical imaging sources \cite{yang2023medmnistv2}. 
It includes diverse modalities such as X-ray, pathology, and organ imaging. 
We use selected subsets to construct few-shot medical target domains, simulating realistic scenarios where only limited annotated medical images are available.

\subsection{Backbone: Latent Diffusion Model}

Stable Diffusion operates in a latent space. 
An image $x$ is encoded into a latent variable $z = \mathcal{E}(x)$ via a pre-trained VAE encoder. 
The diffusion model is trained to predict additive Gaussian noise $\epsilon$ at timestep $t$:

\begin{equation}
\mathcal{L}_{\text{diff}}(\theta) 
= \mathbb{E}_{x,t,\epsilon}
\left[
\|\epsilon - \epsilon_\theta(z_t, t, c)\|_2^2
\right],
\end{equation}

where:
\begin{itemize}
\item $z_t$ is the noisy latent at timestep $t$,
\item $c$ denotes the text conditioning embedding,
\item $\epsilon_\theta$ is the UNet noise predictor.
\end{itemize}

\subsection{Baselines}

\textbf{Zero-shot inference.} 
We directly sample from $p_{\theta_0}(x \mid c)$ without adaptation.

\textbf{Naive few-shot fine-tuning.} 
We optimize all UNet parameters $\theta$ on $\mathcal{D}_t$ using $\mathcal{L}_{\text{diff}}$, allowing unrestricted updates:
\[
\theta = \theta_0 + \Delta\theta, \quad 
\Delta\theta \in \mathbb{R}^{d}.
\]
This serves as a reference for full-capacity adaptation under extreme data scarcity.

\subsection{LoRA-based Parameter-Efficient Adaptation}

Instead of full-parameter updates, we introduce Low-Rank Adaptation (LoRA) into selected attention layers of the UNet. 
For a given weight matrix $W \in \mathbb{R}^{d \times k}$ in self-attention or cross-attention projections, full fine-tuning learns an unrestricted update $\Delta W$. 
In contrast, LoRA parameterizes the update as

\begin{equation}
\Delta W = A B,
\end{equation}

where:
\[
A \in \mathbb{R}^{d \times r}, \quad
B \in \mathbb{R}^{r \times k}, \quad
r \ll \min(d,k).
\]

Thus, the effective weight becomes:
\[
W' = W + A B,
\]
with $\text{rank}(\Delta W) \le r$.

All original backbone parameters are frozen, and only $\{A,B\}$ are optimized. 
The training objective remains the standard diffusion loss $\mathcal{L}_{\text{diff}}$, but the optimization is constrained to a low-dimensional subspace of parameter updates.

\subsection{Algorithmic Procedure}

For each target image $x \in \mathcal{D}_t$:
\begin{enumerate}
    \item Encode $x$ into latent $z = \mathcal{E}(x)$.
    \item Sample timestep $t$ and noise $\epsilon$.
    \item Construct $z_t$ via the forward diffusion process.
    \item Compute $\mathcal{L}_{\text{diff}}$ using the LoRA-augmented UNet.
    \item Update LoRA parameters $\{A,B\}$ via gradient descent.
\end{enumerate}

\subsection{Interpretation as Structural Constraint}

Full fine-tuning permits updates in the full parameter space $\mathbb{R}^{|\theta|}$, which may encourage rapid memorization under extremely small $N$. 
LoRA restricts updates to a low-rank subspace, effectively imposing a structural constraint on how the generation distribution can shift. 
We hypothesize that this reduced update freedom acts as an implicit regularizer, mitigating overfitting while retaining the original generative prior.

\subsection{Evaluation}

To assess few-shot domain adaptation under extreme data scarcity, we evaluate our method along three axes: domain fidelity, generative diversity, and prompt controllability.

\textbf{Domain Fidelity.}
We measure how closely generated samples match the target-domain distribution using two metrics. 
First, we compute Fréchet Inception Distance (FID) between generated samples and real target-domain images. Lower FID indicates closer alignment with the target distribution. 
Second, we compute CLIP-based domain similarity by measuring cosine similarity between generated images and target-domain images in the CLIP image embedding space. Higher similarity suggests stronger domain alignment.

\textbf{Generative Diversity.}
To detect potential mode collapse under extreme few-shot supervision, we compute intra-set LPIPS among generated samples. Higher average LPIPS indicates greater perceptual diversity across outputs.

\textbf{Prompt Controllability.}
To evaluate whether the adapted model remains responsive to textual input, we compute CLIP text–image alignment between generated samples and their conditioning prompts. Higher cosine similarity indicates better preservation of prompt controllability after adaptation.

Together, these metrics provide a multi-dimensional assessment of adaptation quality, capturing distributional shift, diversity preservation, and semantic responsiveness.

\section{Discussion}

\subsection{Novelty}

Our work differs from prior personalization and few-shot diffusion adaptation methods in its focus on structural parameter constraints rather than objective-level regularization. 
While previous approaches either update the full model parameters (e.g., DreamBooth) or introduce additional diversity-preserving losses (e.g., regularized diffusion adaptation), we investigate whether low-rank parameterization alone is sufficient to stabilize extreme few-shot domain shift. 
Instead of proposing a new loss or auxiliary network, we reinterpret LoRA as a structural constraint mechanism and systematically analyze its behavior under severe data scarcity.

\subsection{Advantages}

The proposed method offers several practical and conceptual advantages. 
First, parameter efficiency significantly reduces memory and computational requirements compared to full fine-tuning. 
Second, freezing backbone parameters helps preserve the pre-trained generative prior, mitigating catastrophic drift and prompt collapse. 
Third, the modularity of LoRA enables domain-specific weights to be stored and swapped independently, facilitating scalable multi-domain adaptation.

\subsection{Limitations}

Despite these advantages, several limitations remain. 
Under extremely small and homogeneous datasets, even low-rank updates may still lead to memorization. 
Restricting updates to attention layers may limit expressivity for domains requiring substantial structural changes beyond stylistic shifts. 
Moreover, prompt controllability is partially evaluated using embedding-based similarity metrics, which may not fully capture semantic nuance. 
Future work could explore hybrid approaches combining structural constraints with objective-level regularization.



\bibliography{references}


\end{document}